\documentclass[10pt,a4paper]{article}
\usepackage{cv}
\setdefaultlanguage{russian}
\hypersetup{pdflang={ru}}

\begin{document}
\cvheader{Роман Дугаев}{Инженер-программист · C++ / Python · Сети и backend}

\section*{О себе}
Разрабатываю на C++ и Python протоколы передачи данных, асинхронные сервисы
и системы обработки данных. Занимаюсь переработкой унаследованных библиотек
и созданием новых сервисов: от низкоуровневого сетевого кода до Python API
и инструментов сборки и поставки. Есть опыт в телекоммуникациях, SDR и
разработке аппаратуры на FPGA.

\section*{Опыт работы}
\cventry{DEXOR — IT Systems Engineer}{Февраль 2021 — настоящее время}{Кётен, Германия · Удалённая работа в команде}
Разработка и поддержка сетевых библиотек и backend-сервисов на Python:
реализация, интеграция, тестирование, сборка пакетов и CI/CD.

\cventry{Протоколы передачи данных и транспортные библиотеки}{}{C++, Python, Asio, TCP/UDP, nanobind, CMake, Conan}
\begin{itemize}
  \item Переписал унаследованные библиотеки RSDT и DMTP, общие протокольные,
    интерфейсные и транспортные слои; поддерживал код передачи данных
    «точка — точка» и «точка — многоточка».
  \item Разработал PBTP и RSDT-X: асинхронные C++ библиотеки с параллельными
    TCP-потоками и надёжной доставкой по UDP, корутинами, ограниченными
    буферами и распределением работы между потоками.
  \item Создал Python-привязки и transport-py — общий asyncio API
    с валидацией конфигурации; автоматизировал сборку wheel-пакетов
    с нативным кодом и тестирование.
\end{itemize}

\cventry{DataClone — платформа передачи файлов}{}{Совместная разработка · Python, FastAPI, PostgreSQL, Redis, WebSockets}
\begin{itemize}
  \item Участвовал в разработке backend, агентов, среды передачи данных,
    библиотеки обмена сообщениями, общих контрактов и упаковки Python-клиента:
    backend-py, agent-py, transfer-py, tunnel-py, contracts-py, client-py.
  \item Работал над сканированием и наблюдением за файловой системой,
    версионированием метаданных, управлением процессами передачи,
    пакетной отправкой файлов и отчётами о прогрессе.
  \item Интегрировал нативные транспорты PBTP/RSDT-X с asyncio-сервисами,
    каналами управления WebSocket, ZeroMQ IPC и контрактами Pydantic.
\end{itemize}

\cventry{Aerodata DataCore — платформа обработки данных}{}{Python, FastAPI, SQLModel/SQLAlchemy, PostgreSQL, Temporal, Docker}
\begin{itemize}
  \item Разрабатывал компоненты преобразования, импорта, базы данных
    и backend для авиационных и морских данных: валидацию схем,
    преобразование HSQLDB в SQLite, потоки ZeroMQ и интеграцию импорта в PostgreSQL.
  \item Создал API для рейсов, судов, файлов и задач обработки с аутентификацией
    JWT/Keycloak, ролевым доступом и S3-совместимым хранилищем.
  \item Организовал запуск контейнеров валидации, преобразования и импорта
    через Temporal с повторными попытками, отменой и отслеживанием статуса задач.
\end{itemize}

\newpage
\section*{Технические навыки}
\cvskill{Основные языки}{C++17–23, C++26 в новых транспортных библиотеках;
  Python 3.11+, asyncio, типизация и Pydantic.}
\cvskill{Сети}{TCP/UDP, Asio/Boost.Asio, корутины, многопоточность,
  POSIX/Linux, WebSockets, ZeroMQ, Protobuf и gRPC.}
\cvskill{Backend и данные}{FastAPI, SQLAlchemy/SQLModel, PostgreSQL,
  SQLite, Redis, Temporal, S3-совместимые хранилища; ранее — InfluxDB.}
\cvskill{Сборка и поставка}{CMake, Conan, nanobind, Python-пакеты,
  PyInstaller, Docker/Compose, Git, GitLab CI и GitHub Actions.}
\cvskill{Тестирование и инструменты}{GoogleTest, pytest, Linux, Bash,
  Jupyter, Plotly и техническая документация.}
\cvskill{Интерфейсы}{Практический опыт React и TypeScript для несложных
  веб-интерфейсов; PySide и QML для настольных инструментов.}
\cvskill{Rust — начальный уровень}{Изучаю язык в собственных экспериментах
  с Tokio, UDP, сериализацией пакетов и асинхронными тестами.}
\cvskill{AI в разработке}{Использую AI-агентов для работы с кодом
  и AI-чаты как часть рабочего процесса.}

\section*{Предыдущие инженерные проекты}
\cventry{Производственные инструменты и веб-приложения}{DEXOR}{Python, PySide/QML, React/TypeScript}
\begin{itemize}
  \item Разработал инструмент тестирования авиационного мультимедийного
    оборудования с графическим и консольным интерфейсами на Python и QML.
  \item Работал над приложениями для управления фермой и мониторинга солнечных
    панелей, включая несложные React-интерфейсы и Python-сервисы.
\end{itemize}

\cventry{SDR-плата на FPGA}{Июль 2020 — апрель 2022}{СибГУТИ / Hochschule Anhalt · Altium Designer, SystemVerilog, Quartus, Questa}
\begin{itemize}
  \item Спроектировал шестислойную SDR-плату с FPGA, радиочастотным трактом,
    фильтрами, Ethernet, USB и GPIO; изготовил и протестировал прототип.
  \item Рассчитал радиочастотные фильтры, разработал цифровую часть
    и проверил базовую функциональность прототипа.
\end{itemize}

\cventry{Беспроводной протокол на SDR}{Май 2019 — июль 2020}{СибГУТИ · C/C++, Verilog, HackRF, Intel Cyclone 10, BeagleBone PRU}
Разработал беспроводной протокол, доработал прошивку HackRF и реализовал
синхронизацию несущей, битов и кадров с обработкой данных на FPGA и PRU.

\section*{Образование}
\cventry{Сибирский государственный университет}{2013 — 2020}{телекоммуникаций и информатики (СибГУТИ) · Новосибирск, Россия}
Бакалавриат и магистратура: телекоммуникации, сети и системы коммутации.
Также занимался исследовательской работой в аспирантуре.

\section*{Языки и интересы}
\cvskill{Языки}{Русский — родной; английский — B1.}
\cvskill{Интересы}{Linux, электроника и SDR, Home Assistant и домашняя автоматизация.}

\end{document}
